% ----------------------------------------------------------
\chapter{Fundamentação Teórica}
% ----------------------------------------------------------

A modelagem de dados é a atividade que antecede a construção de um banco de dados e consiste em representar, de forma abstrata, os objetos de um domínio e as associações existentes entre eles. Conforme \citeonline{Fundamentos2017}, essa atividade é organizada em três níveis: o conceitual, que descreve o domínio sem vinculação a uma tecnologia específica e tem no diagrama ER sua principal forma de representação; o lógico, que traduz esse modelo para as estruturas de um determinado tipo de banco de dados; e o físico, que define como os dados serão efetivamente armazenados. O diagrama ER ocupa, portanto, uma posição central no projeto de um sistema, pois é a partir dele que as demais decisões sobre o armazenamento dos dados são derivadas.

Além de orientar a construção do banco de dados, o diagrama ER cumpre um papel de comunicação: é sobre ele que analistas, desenvolvedores e demais envolvidos discutem o domínio e constroem um entendimento comum a respeito dele. Essa natureza compartilhada faz da modelagem uma atividade essencialmente coletiva e explica por que a possibilidade de várias pessoas trabalharem sobre um mesmo modelo é uma demanda recorrente nas ferramentas de apoio ao projeto de bancos de dados.

Os sistemas que dão suporte a esse tipo de trabalho coletivo são objeto de estudo da área de \textit{groupware}. Segundo \citeonline{Fuks2002}, tais sistemas se apoiam em três dimensões — comunicação, coordenação e cooperação — que precisam ser contempladas de maneira articulada: não basta permitir que os participantes atuem sobre um artefato comum, é necessário também coordenar essa atuação para que as contribuições não se anulem. No caso da edição simultânea, a coordenação se materializa em mecanismos de controle de concorrência, responsáveis por impedir que alterações realizadas ao mesmo tempo sobre um mesmo elemento conduzam o artefato a um estado inconsistente \cite{Filippo2005}.

Quando essas dimensões são atendidas, o ganho para a equipe é expressivo: a colaboração em tempo real encurta o ciclo de troca de ideias, oferece visibilidade imediata sobre a contribuição de cada participante e reduz o retrabalho decorrente de versões divergentes de um mesmo artefato \cite{Inoue2014,Hsu2021}. São esses conceitos que orientam a análise das ferramentas de modelagem de diagramas ER apresentada nas seções seguintes.

\section[Análise dos Resultados]{Análise dos Resultados}

Esta seção analisa comparativamente três ferramentas de modelagem de diagramas ER disponíveis no mercado (brModeloWeb, SqlDBM e Vertabelo) com o objetivo de identificar seus pontos fortes e fracos em relação à modelagem colaborativa.

\subsection[brModeloWeb]{brModeloWeb}

O brModeloWeb destaca-se como uma ferramenta de modelagem de dados \textit{online}, acessível e de código aberto, amplamente utilizada em ambientes acadêmicos e de ensino. Sua interface intuitiva e a capacidade de gerar \textit{scripts} SQL a partir dos diagramas ER o tornam uma opção atraente para iniciantes e para projetos de menor escala. No entanto, em termos de funcionalidades colaborativas, o brModeloWeb apresenta limitações significativas. Embora permita o compartilhamento de diagramas e a adição de comentários, não oferece suporte à edição simultânea por múltiplos usuários, controle de versão integrado ou recursos avançados de gerenciamento de conflitos. Essa falta de recursos colaborativos mais robustos pode restringir sua aplicabilidade em projetos de grande porte ou em equipes geograficamente distribuídas. Sendo assim, a colaboração se resume ao compartilhamento de arquivos, o que não permite a rastreabilidade das alterações e dificulta a colaboração em tempo real.

\subsection[SqlDBM]{SqlDBM}

O SqlDBM se apresenta como uma solução de modelagem de banco de dados \textit{online} mais abrangente, com foco na colaboração e na integração com diferentes SGBDs. Uma de suas principais vantagens é o suporte à edição simultânea de diagramas ER por múltiplos usuários, o que facilita a colaboração em tempo real e agiliza o processo de modelagem. Adicionalmente, o SqlDBM oferece controle de versão integrado, permitindo rastrear as alterações realizadas ao longo do tempo e reverter para versões anteriores do modelo, se necessário. A integração com o sistema de controle de versão \textit{git} permite que as equipes de desenvolvimento gerenciem seus modelos de dados da mesma forma que gerenciam o código-fonte da aplicação.

\subsection[Vertabelo]{Vertabelo}

O Vertabelo se posiciona como uma ferramenta de modelagem de dados online que engloba desde a modelagem conceitual (diagramas ER) até a modelagem lógica e física de bancos de dados. Similarmente ao SqlDBM, o Vertabelo oferece suporte à edição simultânea de diagramas, permitindo que múltiplos usuários colaborem em tempo real na criação e na modificação dos modelos de dados. A ferramenta também oferece funcionalidades de validação de modelos, garantindo a consistência e a integridade dos diagramas ER. Além disso, o Vertabelo se destaca pela capacidade de gerar documentação completa dos modelos de dados, facilitando a comunicação e o entendimento entre os membros da equipe de desenvolvimento.

\section[Análise Comparativa]{Análise Comparativa}

A análise das ferramentas brModeloWeb, SqlDBM e Vertabelo revela diferentes abordagens para a modelagem ER colaborativa. Enquanto o brModeloWeb oferece uma solução mais simples e focada em modelagem básica, o SqlDBM e o Vertabelo apresentam funcionalidades mais avançadas de colaboração, como edição em tempo real e controle de versão.

\begin{table}[H]
    \centering
    \renewcommand{\arraystretch}{1.5}
        \begin{tabulary}{\textwidth}{|L|C|C|C|}
        \hline
        \textbf{Funcionalidade} & \textbf{brModeloWeb} & \textbf{SqlDBM} & \textbf{Vertabelo} \\
        \hline
        Colaboração em Tempo Real & Não & Sim & Sim \\
        \hline
        Controle de Versão & Não & Sim & Sim \\
        \hline
        Geração de SQL & Sim & Sim & Sim \\
        \hline
        Todas funcionalidades no pacote grátis & Sim & Não & Não \\
        \hline
        \end{tabulary}
    \caption{Comparativo das Ferramentas}
    \label{tabela:comparativo-ferramentas}
\end{table}
